\chapter{Dependencias lógicas y certificados de consistencia}

\section{Cadena causal de los certificados: estados finitos \(\to\) restricciones \(\to\) supervivencia \(\to\) sección infinita \(\to\) realizaciones}

La cadena lógica de la integración no parte de una conclusión arquimediana.
Se organiza en el orden siguiente:
\[
\begin{tikzpicture}[baseline=(current bounding box.center),>=Latex,
 node distance=7mm and 8mm,
 n/.style={rounded corners,draw=hmtblue,fill=hmtlightblue,
 minimum height=8mm,align=center,font=\small}]
 \node[n] (finite) {estados finitos\\tipados};
 \node[n,right=of finite] (maps) {restricciones\\compatibles};
 \node[n,right=of maps] (surv) {supervivencia\\y poda};
 \node[n,right=of surv] (lim) {sección\\infinita};
 \node[n,right=of lim] (read) {lectores\\y realizaciones};
 \draw[->,thick] (finite)--(maps);
 \draw[->,thick] (maps)--(surv);
 \draw[->,thick] (surv)--(lim);
 \draw[->,thick] (lim)--(read);
\end{tikzpicture}
\]
Un certificado de cifras verifica una restricción de la última flecha; no
crea retrospectivamente los estados ni sustituye el teorema de existencia
del límite.

\section{Identidades aritméticas exactas de los certificados finitos}

Los censos que acompañan al certificado de filtración nonádica satisfacen
\[
 396\cdot6=2376,
 \qquad
 43\cdot9=387,
 \qquad
 432+4\cdot144=1008.
\]
La primera igualdad relaciona eventos elementales de frontera y trits
emitidos por canal; la segunda cuenta las parejas de fase \((K,K+9)\); la
tercera reconstruye la multiplicidad total de las cinco regiones de
\(\pi\).

Los dos papeles regionales que no pueden permutarse son
\[
\begin{array}{rcl}
U_\pi^\perp&=&(501,614,498,169,272,272),\\
E_{\rm sig}&=&555555,\qquad \text{multiplicidad }432,\\[.35em]
U_\pi^{--}&=&(870,923,810,418,674,521),\\
E_{\rm sig}&=&933717,\qquad \text{multiplicidad }144.
\end{array}
\]

\section{Existencia y unicidad de la rama infinita en un árbol finitamente ramificado de prefijos supervivientes}

\begin{lemma}[Árbol finitamente ramificado]
Sea \(T\) un árbol con niveles finitos no vacíos y tal que todo nodo
superviviente tiene un descendiente superviviente. Entonces existe una rama
infinita. Si, además, cada nodo de la subtorre seleccionada tiene exactamente
un descendiente compatible, la rama es única.
\end{lemma}

\begin{proof}
La existencia es el lema de König aplicado al árbol finitamente ramificado.
Para la unicidad, dos ramas distintas tendrían un primer nivel en el que
divergen, dando dos descendientes compatibles del mismo prefijo, contra la
hipótesis.
\end{proof}

\begin{controlbox}
Una ejecución a profundidad \(N\) sólo prueba lo ocurrido en el árbol
truncado. La afirmación infinita requiere la ley uniforme que preserva no
vaciedad y compatibilidad para todo nivel. El cuerpo separa expresamente
ambos estatutos y, para la filtración nonádica, conserva su selector uniforme
certificado.
\end{controlbox}

\section{Clausura y potencia plena}

Para \(A\subseteq X=\varprojlim X_n\), se ha probado
\[
 \overline A=\bigcap_n p_n^{-1}(p_n(A)).
\]
El conjunto de sucesiones binarias finalmente nulas es denso y propio en
\(2^{\mathbb N}\), pero posee todas las sombras finitas. Por ello cualquier
algoritmo que reciba sólo \((p_n(A))_n\) no puede distinguir \(A\) de su
clausura. Este ejemplo bloquea la falsa identidad
\[
 \varprojlim\mathcal P(X_n)=\mathcal P(\varprojlim X_n).
\]

\section{Compatibilidad tipada de los operadores}

La diferencia entre la inclusión unital y la esquina marcada se verifica
directamente:
\[
\begin{aligned}
 \tau_{d+1}(a\otimes I_9)
 &=\frac{\operatorname{Tr}(a)\operatorname{Tr}(I_9)}{9^{d+1}}
 =\tau_d(a),\\
 \tau_{d+1}(a\otimes p_j)
 &=\frac{\operatorname{Tr}(a)}{9^{d+1}}
 =\frac1{9}\tau_d(a).
\end{aligned}
\]
La primera rama es unital y tracial; la segunda conserva una esquina marcada
y no puede usarse para inferir el factor \(II_1\).

\section{Anulación del coste de Landauer en la inversión reversible del estado}

Para una variable completa \(X\) y una proyección \(Y=q(X)\),
\[
 H(X)=H(Y)+H(X\mid Y).
\]
Si \(q\) es biyectiva, la fibra es unitaria y \(H(X\mid Y)=0\). Si \(q\)
es un reinicio del TRIT uniforme, la información descartada es \(\log3\) y
el coste mínimo es \(k_B\Theta\log3\). Para la extensión
\(\widehat E_N\) a \(L^1\) de una esperanza condicional normal
traza-preservante \(E_N\), y bajo las hipótesis de finitud declaradas,
\[
 S_\tau(\widehat E_Nh)-S_\tau(h)
 =D_\tau(h\Vert \widehat E_Nh)\ge0.
\]
Las dos fórmulas localizan el mismo fenómeno en dominios diferentes: el
coste pertenece a la contracción de una fibra, no a la evolución reversible
del estado completo.

\section{Obstrucción Gödel–Turing a un decisor uniforme de multiplicidades proyectivas}

Sea \(T_e\) la máquina de índice \(e\). En el nivel \(n\) se define
\[
 X_n^e=\{b_n\}\cup
 \{c_n:T_e\text{ no se detuvo antes de }n\}.
\]
Cada \(X_n^e\) es computable. La fibra global tiene dos ramas si \(T_e\) no
se detiene y una si se detiene. Por tanto
\[
 U_e:\quad \#q_\infty^{-1}(v)=1
 \quad\Longleftrightarrow\quad T_e\downarrow.
\]
Un decisor uniforme para \(U_e\) decidiría \(\HALT\). Una teoría
recursivamente axiomatizable, correcta y completa para todos los \(U_e\)
permitiría decidirlos enumerando en paralelo pruebas de \(U_e\) y de su
negación. Éste es el núcleo del \emph{Teorema HMT de incompletitud uniforme
para multiplicidades proyectivas}; su demostración emplea una reducción
Gödel--Turing. La obstrucción uniforme no niega los selectores particulares
construidos sobre imágenes HMT calibradas.

\begin{controlbox}[title={Alcance del control de incompletitud uniforme}]
La reducción Gödel--Turing establece la obstrucción uniforme para las
multiplicidades proyectivas. No niega los selectores particulares construidos
sobre imágenes HMT calibradas ni convierte una obstrucción de decisión global
en una indeterminación de cada fibra concreta.
\end{controlbox}
